A supply chain coordinator opens the week with a list of nodes and a finite amount of
attention. Somewhere on that list is the supplier that will cause next month's line
stoppage. Nothing on the screen says which one it is.

That is the practical shape of the problem this mission addressed. Modern supply chains are
monitored through indicators that each measure one dimension at a time, while the decision
the coordinator actually faces is multi-dimensional: cost against deadline against supplier
reliability against, increasingly, carbon. When the tooling cannot combine those dimensions,
the coordinator combines them by judgment. Judgment is then the real decision layer of the
system, and it is the layer nobody measures.

The consequences are not evenly distributed. Declaring more urgency than a situation
warrants consumes shared transport and production capacity that another flow needed.
Declaring less leaves a risk that nobody is watching, and by the time it becomes visible the
window to act on it has closed. The second error is the more expensive, and it is the harder
one to detect, precisely because an under-declared risk generates no signal.

This mission was carried out at the Direction de l'Innovation of ALTEN, within the Smart
Green Supply Chain research programme, between 13 April and 16 October 2026. Its subject was
the layer that sits before any planning decision: the urgency signal itself. The question
driving the work was whether the gap between what an operator declares and what the
operational state actually supports can be measured well enough, inside a working digital
twin, to be acted upon.

The answer this report reaches is a qualified no, and the value of the work lies in how
precisely that no can be stated. A complete system was delivered and is operational. The
signal it produces was then tested against recorded outcomes rather than asserted, and the
test failed in a way that identifies exactly which component is responsible and what would
repair it. A second measurement, unplanned at the outset, showed that displaying the
system's own forecast to the person making the declaration changed that declaration, and
changed it in the wrong direction at the worst possible moment. Both findings are more
useful to ALTEN than a favourable score would have been.

The report is organised as the mission itself was. Section~\ref{sec:entreprise} presents
ALTEN, the Direction de l'Innovation and the research programme this work sits inside,
including the company's corporate social responsibility policy and how the mission relates
to it. Section~\ref{sec:problematique} states the problem that justified opening the
subject. Section~\ref{sec:mission} gives the assignment, its objectives, its indicators and
its deliverables. Section~\ref{sec:organisation} describes how the project was run: who was
involved, what resources were available, which working method was applied, how the six
months were phased, and which risks were carried. Section~\ref{sec:etatdelart} reviews the
existing methods and identifies what none of them provides.
Section~\ref{sec:realisation} reports the work carried out and what it produced, including
the problems met along the way and how they were resolved. Section~\ref{sec:conclusion}
assesses the results against the objectives, sets out what remains to be done, and closes
with a personal assessment against the ESTIA engineering competency framework.
